\subsection*{\texorpdfstring{Divisibility in $\mathbb{Z}$}{Divisibility in Z}}

\begin{definitionbox}{Definition}
\begin{definition}[Divisibility in $\mathbb{Z}$]
\label{def:divisibility-in-z}
For $a,b\in\mathbb{Z}$, define
\[
a\mid b
\quad:\Longleftrightarrow\quad
\exists c\in\mathbb{Z}\;(b=a\cdot c).
\]
\end{definition}
\end{definitionbox}

\begin{remark*}[Standard quantified statement]
\[
\forall a,b\in\mathbb{Z},\quad
a\mid b \Longleftrightarrow \exists c\in\mathbb{Z}\,(b=a\cdot c).
\]
\end{remark*}

\begin{remark*}[Predicate reading]
\[
\operatorname{Divides}(a,b,\mathbb{Z},\cdot_{\mathbb{Z}}).
\]
\end{remark*}

\begin{remark*}[Interpretation]
Divisibility records the existence of an integer multiplier from the divisor to
the dividend.
\end{remark*}

\begin{dependencies}
\begin{itemize}
  \item \hyperref[def:integers]{Integers}
  \item \hyperref[thm:multiplication-on-z-well-defined]{Multiplication on $\mathbb{Z}$ Is Well-Defined}
  \item \hyperref[def:existential-q]{Existential Formula}
\end{itemize}
\end{dependencies}

\subsection*{Basic Divisibility Laws}

\begin{propositionbox}{Proposition}
\begin{proposition}[Divisibility Is Reflexive on $\mathbb{Z}$]
\label{prop:divisibility-reflexive-on-z}
For every $a\in\mathbb{Z}$,
\[
a\mid a.
\]
\hyperref[prf:divisibility-reflexive-on-z]{\textit{Go to proof.}}
\end{proposition}
\end{propositionbox}

\begin{remark*}[Standard quantified statement]
\[
\forall a\in\mathbb{Z},\quad a\mid a.
\]
\end{remark*}

\begin{remark*}[Predicate reading]
\[
\operatorname{Reflexive}(\mid,\mathbb{Z}).
\]
\end{remark*}

\begin{remark*}[Interpretation]
Every integer divides itself by multiplication with one.
\end{remark*}

\begin{dependencies}
\begin{itemize}
  \item \hyperref[def:divisibility-in-z]{Divisibility in $\mathbb{Z}$}
  \item \hyperref[def:one-in-z]{One in $\mathbb{Z}$}
  \item \hyperref[thm:z-multiplicative-commutative-monoid]{$\mathbb{Z}$ Multiplicative Commutative Monoid}
\end{itemize}
\end{dependencies}

\begin{propositionbox}{Proposition}
\begin{proposition}[Divisibility Is Transitive on $\mathbb{Z}$]
\label{prop:divisibility-transitive-on-z}
For all $a,b,c\in\mathbb{Z}$,
\[
a\mid b\land b\mid c \Longrightarrow a\mid c.
\]
\hyperref[prf:divisibility-transitive-on-z]{\textit{Go to proof.}}
\end{proposition}
\end{propositionbox}

\begin{remark*}[Standard quantified statement]
\[
\forall a,b,c\in\mathbb{Z},\quad
a\mid b\land b\mid c \Longrightarrow a\mid c.
\]
\end{remark*}

\begin{remark*}[Predicate reading]
\[
\operatorname{Transitive}(\mid,\mathbb{Z}).
\]
\end{remark*}

\begin{remark*}[Interpretation]
Divisibility chains compose through multiplication.
\end{remark*}

\begin{dependencies}
\begin{itemize}
  \item \hyperref[def:divisibility-in-z]{Divisibility in $\mathbb{Z}$}
  \item \hyperref[thm:multiplication-on-z-associative]{Multiplication on $\mathbb{Z}$ Is Associative}
\end{itemize}
\end{dependencies}

\begin{propositionbox}{Proposition}
\begin{proposition}[One and Negative One Divide Every Integer]
\label{prop:units-divide-everything-on-z}
For every $a\in\mathbb{Z}$,
\[
1\mid a
\qquad\text{and}\qquad
(-1)\mid a.
\]
\hyperref[prf:units-divide-everything-on-z]{\textit{Go to proof.}}
\end{proposition}
\end{propositionbox}

\begin{remark*}[Standard quantified statement]
\[
\forall a\in\mathbb{Z},\quad 1_{\mathbb{Z}}\mid a\land (-1_{\mathbb{Z}})\mid a.
\]
\end{remark*}

\begin{remark*}[Predicate reading]
\[
\operatorname{UnitsDivideAll}(\mathbb{Z},\cdot_{\mathbb{Z}},1_{\mathbb{Z}}).
\]
\end{remark*}

\begin{remark*}[Interpretation]
The two integer units divide every integer.
\end{remark*}

\begin{dependencies}
\begin{itemize}
  \item \hyperref[def:divisibility-in-z]{Divisibility in $\mathbb{Z}$}
  \item \hyperref[def:one-in-z]{One in $\mathbb{Z}$}
  \item \hyperref[cor:negative-one-rule-on-z]{Negative One Rule on $\mathbb{Z}$}
\end{itemize}
\end{dependencies}

\begin{propositionbox}{Proposition}
\begin{proposition}[Every Integer Divides Zero]
\label{prop:everything-divides-zero-on-z}
For every $a\in\mathbb{Z}$,
\[
a\mid 0.
\]
\hyperref[prf:everything-divides-zero-on-z]{\textit{Go to proof.}}
\end{proposition}
\end{propositionbox}

\begin{remark*}[Standard quantified statement]
\[
\forall a\in\mathbb{Z},\quad a\mid0_{\mathbb{Z}}.
\]
\end{remark*}

\begin{remark*}[Predicate reading]
\[
\operatorname{DividesZero}(a,\mathbb{Z},\cdot_{\mathbb{Z}},0_{\mathbb{Z}}).
\]
\end{remark*}

\begin{remark*}[Interpretation]
Every integer divides zero because zero is absorbing for multiplication.
\end{remark*}

\begin{dependencies}
\begin{itemize}
  \item \hyperref[def:divisibility-in-z]{Divisibility in $\mathbb{Z}$}
  \item \hyperref[cor:zero-absorption-on-z]{Zero Absorption on $\mathbb{Z}$}
\end{itemize}
\end{dependencies}

\begin{propositionbox}{Proposition}
\begin{proposition}[Divisibility Is Compatible with Negation on the Divisor]
\label{prop:divisibility-compatible-with-negation-divisor-on-z}
For all $a,b\in\mathbb{Z}$,
\[
a\mid b \Longleftrightarrow (-a)\mid b.
\]
\hyperref[prf:divisibility-compatible-with-negation-divisor-on-z]{\textit{Go to proof.}}
\end{proposition}
\end{propositionbox}

\begin{remark*}[Standard quantified statement]
\[
\forall a,b\in\mathbb{Z},\quad
a\mid b \Longleftrightarrow (-a)\mid b.
\]
\end{remark*}

\begin{remark*}[Predicate reading]
\[
\operatorname{DivisibilityInvariantUnderNegatedDivisor}(\mathbb{Z}).
\]
\end{remark*}

\begin{remark*}[Interpretation]
Changing the sign of the divisor does not change divisibility.
\end{remark*}

\begin{dependencies}
\begin{itemize}
  \item \hyperref[def:divisibility-in-z]{Divisibility in $\mathbb{Z}$}
  \item \hyperref[cor:sign-rule-on-z]{Sign Rule on $\mathbb{Z}$}
\end{itemize}
\end{dependencies}

\begin{propositionbox}{Proposition}
\begin{proposition}[Divisibility Is Compatible with Negation on the Dividend]
\label{prop:divisibility-compatible-with-negation-dividend-on-z}
For all $a,b\in\mathbb{Z}$,
\[
a\mid b \Longleftrightarrow a\mid(-b).
\]
\hyperref[prf:divisibility-compatible-with-negation-dividend-on-z]{\textit{Go to proof.}}
\end{proposition}
\end{propositionbox}

\begin{remark*}[Standard quantified statement]
\[
\forall a,b\in\mathbb{Z},\quad
a\mid b \Longleftrightarrow a\mid(-b).
\]
\end{remark*}

\begin{remark*}[Predicate reading]
\[
\operatorname{DivisibilityInvariantUnderNegatedDividend}(\mathbb{Z}).
\]
\end{remark*}

\begin{remark*}[Interpretation]
Changing the sign of the dividend does not change divisibility.
\end{remark*}

\begin{dependencies}
\begin{itemize}
  \item \hyperref[def:divisibility-in-z]{Divisibility in $\mathbb{Z}$}
  \item \hyperref[cor:sign-rule-on-z]{Sign Rule on $\mathbb{Z}$}
\end{itemize}
\end{dependencies}

\begin{propositionbox}{Proposition}
\begin{proposition}[Divisibility Is Compatible with Negation on $\mathbb{Z}$]
\label{prop:divisibility-compatible-with-negation-on-z}
For all $a,b\in\mathbb{Z}$,
\[
a\mid b \Longleftrightarrow (-a)\mid b \Longleftrightarrow a\mid(-b).
\]
\hyperref[prf:divisibility-compatible-with-negation-on-z]{\textit{Go to proof.}}
\end{proposition}
\end{propositionbox}

\begin{remark*}[Standard quantified statement]
\[
\forall a,b\in\mathbb{Z},\quad
a\mid b \Longleftrightarrow (-a)\mid b \Longleftrightarrow a\mid(-b).
\]
\end{remark*}

\begin{remark*}[Predicate reading]
\[
\operatorname{DivisibilityInvariantUnderNegation}(\mathbb{Z}).
\]
\end{remark*}

\begin{remark*}[Interpretation]
Divisibility is unchanged by negating either input.
\end{remark*}

\begin{dependencies}
\begin{itemize}
  \item \hyperref[prop:divisibility-compatible-with-negation-divisor-on-z]{Divisibility Is Compatible with Negation on the Divisor}
  \item \hyperref[prop:divisibility-compatible-with-negation-dividend-on-z]{Divisibility Is Compatible with Negation on the Dividend}
\end{itemize}
\end{dependencies}

\begin{propositionbox}{Proposition}
\begin{proposition}[Divisibility Is Compatible with Addition on $\mathbb{Z}$]
\label{prop:divisibility-compatible-with-addition-on-z}
For all $a,b,c\in\mathbb{Z}$,
\[
a\mid b\land a\mid c \Longrightarrow a\mid(b+c).
\]
\hyperref[prf:divisibility-compatible-with-addition-on-z]{\textit{Go to proof.}}
\end{proposition}
\end{propositionbox}

\begin{remark*}[Standard quantified statement]
\[
\forall a,b,c\in\mathbb{Z},\quad
a\mid b\land a\mid c \Longrightarrow a\mid(b+c).
\]
\end{remark*}

\begin{remark*}[Predicate reading]
\[
\operatorname{DivisibilityClosedUnderAddition}(\mathbb{Z},+_{\mathbb{Z}},\cdot_{\mathbb{Z}}).
\]
\end{remark*}

\begin{remark*}[Interpretation]
Common divisibility is preserved under addition of dividends.
\end{remark*}

\begin{dependencies}
\begin{itemize}
  \item \hyperref[def:divisibility-in-z]{Divisibility in $\mathbb{Z}$}

  \item \hyperref[thm:addition-on-z-well-defined]{Addition on $\mathbb{Z}$ Is Well-Defined}
  \item \hyperref[thm:multiplication-distributes-over-addition-on-z]{Multiplication Distributes over Addition on $\mathbb{Z}$}
\end{itemize}
\end{dependencies}

\begin{propositionbox}{Proposition}
\begin{proposition}[Divisibility Respects Linear Combinations on $\mathbb{Z}$]
\label{prop:divisibility-respects-linear-combinations-on-z}
For all $a,b,c,x,y\in\mathbb{Z}$,
\[
a\mid b\land a\mid c \Longrightarrow a\mid(b\cdot x+c\cdot y).
\]
\hyperref[prf:divisibility-respects-linear-combinations-on-z]{\textit{Go to proof.}}
\end{proposition}
\end{propositionbox}

\begin{remark*}[Standard quantified statement]
\[
\forall a,b,c,x,y\in\mathbb{Z},\quad
a\mid b\land a\mid c \Longrightarrow a\mid(b\cdot x+c\cdot y).
\]
\end{remark*}

\begin{remark*}[Predicate reading]
\[
\operatorname{DivisibilityClosedUnderLinearCombinations}(\mathbb{Z},+_{\mathbb{Z}},\cdot_{\mathbb{Z}}).
\]
\end{remark*}

\begin{remark*}[Interpretation]
A common divisor divides every integer linear combination of the two dividends.
\end{remark*}

\begin{dependencies}
\begin{itemize}
  \item \hyperref[def:divisibility-in-z]{Divisibility in $\mathbb{Z}$}
  \item \hyperref[prop:divisibility-compatible-with-addition-on-z]{Divisibility Is Compatible with Addition on $\mathbb{Z}$}
  \item \hyperref[thm:multiplication-on-z-associative]{Multiplication on $\mathbb{Z}$ Is Associative}
  \item \hyperref[thm:multiplication-distributes-over-addition-on-z]{Multiplication Distributes over Addition on $\mathbb{Z}$}
\end{itemize}
\end{dependencies}

\subsection*{\texorpdfstring{Units in $\mathbb{Z}$}{Units in Z}}

\begin{definitionbox}{Definition}
\begin{definition}[Unit in $\mathbb{Z}$]
\label{def:unit-in-z}
An integer $u\in\mathbb{Z}$ is a \emph{unit} if
\[
\exists v\in\mathbb{Z}\;(u\cdot v=1).
\]
\end{definition}
\end{definitionbox}

\begin{remark*}[Standard quantified statement]
\[
\operatorname{Unit}(u)
\Longleftrightarrow
u\in\mathbb{Z}\land\exists v\in\mathbb{Z}\,(u\cdot v=1_{\mathbb{Z}}).
\]
\end{remark*}

\begin{remark*}[Predicate reading]
\[
\operatorname{Unit}(u,\mathbb{Z},\cdot_{\mathbb{Z}},1_{\mathbb{Z}}).
\]
\end{remark*}

\begin{remark*}[Interpretation]
A unit is an integer with a multiplicative inverse.
\end{remark*}

\begin{dependencies}
\begin{itemize}
  \item \hyperref[def:integers]{Integers}
  \item \hyperref[def:one-in-z]{One in $\mathbb{Z}$}
  \item \hyperref[thm:multiplication-on-z-well-defined]{Multiplication on $\mathbb{Z}$ Is Well-Defined}
  \item \hyperref[def:existential-q]{Existential Formula}
\end{itemize}
\end{dependencies}

\begin{theorembox}{Theorem}
\begin{theorem}[The Units of $\mathbb{Z}$ Are Exactly $\pm1$]
\label{thm:units-of-z-are-pm-one}
An integer $u$ is a unit if and only if
\[
u=1\quad\text{or}\quad u=-1.
\]
\hyperref[prf:units-of-z-are-pm-one]{\textit{Go to proof.}}
\end{theorem}
\end{theorembox}

\begin{remark*}[Standard quantified statement]
\[
\forall u\in\mathbb{Z},\quad
\operatorname{Unit}(u,\mathbb{Z},\cdot_{\mathbb{Z}},1_{\mathbb{Z}})
\Longleftrightarrow
u=1_{\mathbb{Z}}\lor u=-1_{\mathbb{Z}}.
\]
\end{remark*}

\begin{remark*}[Predicate reading]
\[
\operatorname{UnitsExactly}(\mathbb{Z},\{1_{\mathbb{Z}},-1_{\mathbb{Z}}\}).
\]
\end{remark*}

\begin{remark*}[Interpretation]
The only invertible integers are one and negative one.
\end{remark*}

\begin{dependencies}
\begin{itemize}
  \item 
\hyperref[def:unit-in-z]{Unit in $\mathbb{Z}$}
  \item \hyperref[thm:z-integral-domain]{$\mathbb{Z}$ Is an Integral Domain}
  \item \hyperref[thm:trichotomy-on-z]{Trichotomy on $\mathbb{Z}$}
  \item \hyperref[cor:negative-one-rule-on-z]{Negative One Rule on $\mathbb{Z}$}
\end{itemize}
\end{dependencies}

\begin{definitionbox}{Definition}
\begin{definition}[Associates in $\mathbb{Z}$]
\label{def:associates-in-z}
Integers $a,b\in\mathbb{Z}$ are \emph{associates} if $a=u\cdot b$ for some
unit $u\in\mathbb{Z}$.
\end{definition}
\end{definitionbox}

\begin{remark*}[Standard quantified statement]
\[
\operatorname{Associates}(a,b)
\Longleftrightarrow
\exists u\in\mathbb{Z},\,
\operatorname{Unit}(u,\mathbb{Z},\cdot_{\mathbb{Z}},1_{\mathbb{Z}})
\land a=u\cdot b.
\]
\end{remark*}

\begin{remark*}[Predicate reading]
\[
\operatorname{Associates}(a,b,\mathbb{Z},\cdot_{\mathbb{Z}},1_{\mathbb{Z}}).
\]
\end{remark*}

\begin{remark*}[Interpretation]
Associate integers differ by multiplication by a unit.
\end{remark*}

\begin{dependencies}
\begin{itemize}
  \item \hyperref[def:unit-in-z]{Unit in $\mathbb{Z}$}
  \item \hyperref[thm:multiplication-on-z-well-defined]{Multiplication on $\mathbb{Z}$ Is Well-Defined}
\end{itemize}
\end{dependencies}

\subsection*{Division Algorithm}

\begin{theorembox}{Theorem}
\begin{theorem}[Division Algorithm on $\mathbb{Z}$]
\label{thm:division-algorithm-on-z}
If $a,b\in\mathbb{Z}$ and $b\ne0$, then there exist unique
$q,r\in\mathbb{Z}$ such that
\[
a=bq+r
\qquad\text{and}\qquad
0\le r<|b|.
\]
\hyperref[prf:division-algorithm-on-z]{\textit{Go to proof.}}
\end{theorem}
\end{theorembox}

\begin{remark*}[Standard quantified statement]
\[
\forall a,b\in\mathbb{Z},\ b\ne0_{\mathbb{Z}}
\Longrightarrow
\exists! q,r\in\mathbb{Z},\,
a=bq+r\land 0\le r<|b|.
\]
\end{remark*}

\begin{remark*}[Predicate reading]
\[
\operatorname{DivisionAlgorithm}(\mathbb{Z},+_{\mathbb{Z}},\cdot_{\mathbb{Z}},\le_{\mathbb{Z}},|\cdot|).
\]
\end{remark*}

\begin{remark*}[Interpretation]
Division by a nonzero integer has unique quotient and remainder with bounded
nonnegative remainder.
\end{remark*}

\begin{dependencies}
\begin{itemize}
  \item \hyperref[thm:well-ordering-principle]{Well-Ordering Principle}
  \item \hyperref[cor:multiplicative-cancellation-on-z]{Multiplicative Cancellation on $\mathbb{Z}$}
  \item \hyperref[thm:z-is-not-well-ordered]{$\mathbb{Z}$ Is Not Well-Ordered}
\end{itemize}
\end{dependencies}

\subsection*{Greatest Common Divisor}

\begin{definitionbox}{Definition}
\begin{definition}[Greatest Common Divisor in $\mathbb{Z}$]
\label{def:greatest-common-divisor-in-z}
A greatest common divisor $\gcd(a,b)$ is the positive common divisor of
$a$ and $b$ divisible by every other common divisor.
\end{definition}
\end{definitionbox}

\begin{remark*}[Standard quantified statement]
\[
d=\gcd(a,b)
\Longleftrightarrow
0<d\land d\mid a\land d\mid b
\land
\forall e\in\mathbb{Z},\ (e\mid a\land e\mid b)\Longrightarrow e\mid d.
\]
\end{remark*}

\begin{remark*}[Predicate reading]
\[
\operatorname{GreatestCommonDivisor}(d,a,b,\mathbb{Z},\mid,<_{\mathbb{Z}}).
\]
\end{remark*}

\begin{remark*}[Interpretation]
A greatest common divisor is the positive common divisor that absorbs all
other common divisors.
\end{remark*}

\begin{dependencies}
\begin{itemize}
  \item \hyperref[def:divisibility-in-z]{Divisibility in $\mathbb{Z}$}
  \item \hyperref[def:positive-integer]{Positive Integer}
  \item \hyperref[def:order-on-z]{Order on $\mathbb{Z}$}
\end{itemize}
\end{dependencies}

\subsection*{Bezout Identity}

\begin{theorembox}{Theorem}
\begin{theorem}[Bezout Identity on $\mathbb{Z}$]
\label{thm:bezout-identity-on-z}
For $a,b\in\mathbb{Z}$, there exist $x,y\in\mathbb{Z}$ such that
\[
\gcd(a,b)=a\cdot x+b\cdot y.
\]
\hyperref[prf:bezout-identity-on-z]{\textit{Go to proof.}}
\end{theorem}
\end{theorembox}

\begin{remark*}[Standard quantified statement]
\[
\forall a,b\in\mathbb{Z}\,\exists x,y\in\mathbb{Z},\quad
\gcd(a,b)=a\cdot x+b\cdot y.
\]
\end{remark*}

\begin{remark*}[Predicate reading]
\[
\operatorname{BezoutIdentity}(\mathbb{Z},+_{\mathbb{Z}},\cdot_{\mathbb{Z}},\gcd).
\]
\end{remark*}

\begin{remark*}[Interpretation]
Every greatest common divisor is an integer linear combination of the inputs.
\end{remark*}

\begin{dependencies}
\begin{itemize}
  \item \hyperref[thm:well-ordering-principle]{Well-Ordering Principle}
\end{itemize}
\end{dependencies}
